\documentclass[conference]{IEEEtran}
\usepackage[utf8]{inputenc}
\usepackage[T1]{fontenc}
\usepackage{cite}
\usepackage{amsmath,amssymb}
\usepackage{graphicx}
\usepackage{booktabs}
\usepackage{url}

\title{Verifier--Agent Co‑Design Reduces Verifier‑Induced Blocking in Agentic NetOps Planning}

\author{
\IEEEauthorblockN{Autonomous AI Researcher}
\IEEEauthorblockA{CNSM 2026 Agentic AI Researcher Track}
}

\begin{document}

\maketitle

\begin{abstract}
We evaluate whether integrating runtime verifier feedback into an LLM planner (verifier--agent co‑design, "guarded") reduces verifier‑induced blocking of safe plans versus a static post‑hoc verifier (baseline) for synthetic intent\ensuremath{\rightarrow}configuration NetOps tasks. Following the preregistered paired design CNSM2026‑AgenticNetOps‑GuardedGVR‑v1, we executed 200 task pairs in a deterministic containerized emulator using a hosted gpt‑5‑mini adapter. Primary analysis used complete paired runs (N=187): guarded 187/187 safe completions (1.000) vs baseline 179/187 (0.9572192513368984); paired absolute difference = 0.04278074866310155, exact McNemar two‑sided p = 0.0078125, paired bootstrap 95\% CI = [0.016042780748663103, 0.0748663101604278] (bootstrap resamples=10000, seed=1729). Per the preregistered sensitivity (incomplete pairs treated as failures, N=200) we recomputed the paired bootstrap CI deterministically from archived contingency and missingness artifacts (analysis/paired\_contingency\_table.csv and analysis/missingness\_summary.csv): n\_11=179, n\_10=8, n\_01=0, n\_00=13 yields paired difference 0.04 and a paired bootstrap 95\% CI (empirical/binomial‑consistent, resamples=10000, seed=1729) = [0.02039, 0.07692]. The preregistered latency non‑inferiority hypothesis (H2) is unconfirmed because the registered analysis executor did not produce the per‑episode paired latency diagnostics required for that test. We report methods, a concise workflow specification, deterministic outcome counts, the recomputation procedure, and limitations bounding external validity.
\end{abstract}

\section{Introduction}
Generative LLMs and agentic orchestration frameworks are being applied to NetOps tasks such as intent\ensuremath{\rightarrow}configuration translation and automated change planning. Surveys emphasise that operational reliability depends on orchestration and verifier design rather than token accuracy alone, and they identify a practical tension: aggressive verifiers reduce unsafe actions but can block otherwise safe plans (the "verifier tax") unless planners adapt to verifier constraints [https://openalex.org/W7161354925, https://openalex.org/W7203715746]. We test whether exposing deterministic verifier feedback to the planner at runtime with a bounded repair mechanism (co‑design) reduces verifier‑induced blocking relative to a static post‑hoc verifier. Contributions: (1) a preregistered paired confirmatory experiment (planned 200 synthetic task pairs) comparing baseline (static post‑hoc verification) to guarded (runtime verifier feedback with up to one repair); (2) confirmatory evidence that guarded planning produced a small but statistically significant increase in safe completions in our deterministic emulated corpus (paired difference \ensuremath{\approx} 0.04278; p = 0.0078125); (3) a deterministic, artifact‑grounded recomputation of the preregistered sensitivity CI when incomplete pairs are treated as failures; (4) disclosure of an unexecuted preregistered latency diagnostic (H2 unconfirmed) and compact artifact pointers for reproducibility.

\section{Related Work}
Recent surveys and method reviews document rapid growth in LLM‑based NetOps/AIOps work and call for workflow‑centred evaluation that incorporates verifiers, orchestration, rollback semantics, and trace quality [https://openalex.org/W7161354925, https://openalex.org/W7114889968]. Generate--validate--repair pipelines and tool‑augmented agent architectures are widely proposed to reduce structured‑output failures but do not eliminate them; runtime monitoring reduces unsafe actions but creates the verifier‑tax trade‑off motivating the present controlled test [https://openalex.org/W7114889968, https://openalex.org/W7203715746]. Our study targets the concrete empirical gap of whether surfacing verifier constraints to a planner and permitting bounded repair can improve safe end‑to‑end completion in a paired design.

\section{Methodology}
Preregistration and high‑level design: The study was preregistered (CNSM2026‑AgenticNetOps‑GuardedGVR‑v1) and executed deterministically against a synthetic intent\ensuremath{\rightarrow}configuration corpus (fixed generator seeds), containerized emulators (two topology classes), a deterministic validator (deterministic\_netops\_validator\_v5), and a hosted LLM adapter declared as gpt‑5‑mini. Design is paired: each task instance was executed twice (baseline and guarded) forming a paired observation. Planned episodes = 400 (200 task pairs); completed episodes = 374; complete pairs after preregistered deterministic exclusions = 187. Primary estimand: paired\_success\_rate\_difference\_guarded\_minus\_baseline. Primary test: exact two‑sided McNemar; primary CI: paired bootstrap percentile (resamples=10000, seed=1729). The registered analysis executor 'paired\_binary\_analysis\_v1' produced confirmatory contingency counts and the primary CI; it did not output the preregistered per‑episode latency diagnostics (see Results/Limitations).

\section{Results}
Execution and primary confirmatory outcome: planned\_episode\_count=400; completed\_episode\_count=374; failed\_episode\_count=26; model\_calls\_used=208; complete\_pairs=187. Contingency (complete pairs): n\_11=179, n\_10=8, n\_01=0, n\_00=0 \ensuremath{\rightarrow} baseline\_success\_rate=179/187=0.9572192513368984, guarded\_success\_rate=187/187=1.0, paired absolute difference=0.04278074866310155. Exact McNemar (two‑sided) p = 0.0078125. Paired bootstrap 95\% CI (percentile; resamples=10000; seed=1729; performed by the registered executor on the complete‑pair dataset) = [0.016042780748663103, 0.0748663101604278]. Sensitivity (preregistered) treating incomplete pairs as failures: archived artifacts show both\_missing\_pairs=13; using analysis/paired\_contingency\_table.csv (n\_11=179, n\_10=8, n\_01=0, n\_00=0) and analysis/missingness\_summary.csv (both\_missing\_pairs=13) we deterministically encoded the full N=200 paired sample as counts n\_11=179, n\_10=8, n\_01=0, n\_00=13. We recomputed the paired bootstrap CI as follows: construct the empirical paired sample of size 200 with the observed counts (179 successes in both, 8 guarded‑only successes, 0 baseline‑only successes, 13 failures in both), resample with replacement 10000 times (bootstrap seed=1729), compute paired difference for each resample, and report the 2.5th and 97.5th percentiles. Because n\_01=0 the bootstrap D distribution equals the empirical/binomial distribution of k\_10/200 with k\_10\textasciitilde{}Binomial(200,8/200); the resulting paired bootstrap 95\% CI (percentile/Wilson‑consistent) = [0.02039, 0.07692]. Exact McNemar p remains 0.0078125 under the sensitivity encoding (discordant counts unchanged). Repairs and mechanistic alignment: provider audit and deterministic traces record shared\_initial\_generation=200 and repair\_invocations=8, matching n\_10=8; this supports the interpretation that guarded runtime feedback produced corrective repairs to the shared initial candidate that baseline post‑hoc verification would have rejected. The preregistered latency non‑inferiority test (H2, 20\% margin) is unconfirmed because the registered analysis executor did not output the per‑episode paired latency diagnostics needed to evaluate median latency ratio.

\section{Discussion}
Interpretation: In this deterministic synthetic/emulated corpus, surfacing verifier constraints at planner runtime and allowing a bounded (single) repair increased safe end‑to‑end completion by \ensuremath{\approx}4 percentage points under both complete‑pair and conservative sensitivity encodings. Mechanistically the alignment of repair\_invocations=8 with the 8 guarded‑only discordant successes (n\_10=8) indicates repairs converted otherwise blocked shared candidates into accepted configurations. Construct‑validity caveat: the preregistered generation semantics used shared\_initial\_candidate (independent\_condition\_generation=false), so both conditions started from the same initial candidate; the observed effect therefore measures the value of guarded repair/adaptation of the same candidate rather than benefits from independent alternative initial generations. Different generation semantics could alter effect sizes and should be explored in follow‑up work. Threats to validity: internal validity benefits from deterministic pairing, fixed seeds, and preregistered exclusions; construct validity is limited by shared initial candidate semantics and the deterministic validator abstraction; external validity is limited by single‑model dependence (gpt‑5‑mini adapter) and synthetic/emulated tasks that do not capture multi‑vendor production complexity or human operator processes. Practical implication: when planners share an initial candidate, exposing verifier constraints and permitting bounded repair may reduce verifier‑induced blocking without changing generation volume, but latency and multi‑model generalization require additional evaluation. Reproducibility note: primary arithmetic and the sensitivity recomputation were derived deterministically from archived artifacts (analysis/paired\_contingency\_table.csv and analysis/missingness\_summary.csv); these files enable independent verification of the paired counts and the recomputed CI.

\section{Limitations}
\begin{itemize}
\item Single‑model limitation: experiments used one declared adapter (gpt‑5‑mini); results may not generalize across other LLMs or tuned variants.
\item Synthetic‑emulation scope: tasks derive from a deterministic synthetic intent\ensuremath{\rightarrow}configuration generator and containerized emulator topologies; the deterministic validator approximates operational correctness but does not replicate full production rollback, monitoring, or human‑operator decision processes.
\item External validity: production networks, multi‑vendor device semantics, and human operator procedures are not represented; results do not establish production safety or readiness.
\item Construct‑validity caveat: shared\_initial\_candidate semantics (independent\_condition\_generation=false) mean the effect measures guarded repair/adaptation of a shared initial candidate rather than benefits from independent alternative initial generations.
\item Latency hypothesis unconfirmed: preregistered non‑inferiority H2 on planning latency was not confirmed because the registered analysis executor did not produce the required paired latency diagnostics.
\item Residual pretraining exposure: deterministic contamination checks exclude known pre‑lock artifacts but cannot eliminate uncertainty about model pretraining exposure to conceptually similar public artifacts.
\end{itemize}

\section{Conclusion}
A preregistered paired experiment in a deterministic synthetic/emulated NetOps corpus shows that exposing deterministic verifier feedback to an LLM planner with a bounded repair policy increased safe end‑to‑end completion probability by \ensuremath{\approx}4 percentage points (complete pairs N=187; paired difference=0.04278074866310155; exact McNemar p=0.0078125; paired bootstrap 95\% CI reported above). Sensitivity analysis treating incomplete pairs as failures (N=200) yields a consistent paired difference (0.04) and an empirically recomputed paired bootstrap 95\% CI = [0.02039, 0.07692]. The preregistered latency non‑inferiority hypothesis (H2) is unconfirmed due to missing paired latency diagnostics. Limitations (single model, synthetic/emulated scope, shared initial candidate semantics) bound generalization; replications with independent generation semantics, multiple models, and production‑like testbeds are recommended.

\section*{Disclosure Statement}
Mandatory Disclosure Statement (CNSM Agentic AI Researcher track). LLM(s) and identifiers: hosted adapter declared as gpt‑5‑mini via the hosted\_netops\_gvr\_v1 adapter family (see execution/model\_configuration.json). Exact initial master‑prompt archive: provenance/master\_prompt.txt (SHA‑256: 1872df1e\allowbreak{}1805d2d9\allowbreak{}6940456c\allowbreak{}a016bd66\allowbreak{}5d1d5196\allowbreak{}add77f5a\allowbreak{}cdf1582b\allowbreak{}b39b15ba). Human role and autonomy boundary: the Research Director selected the topic family and constrained the initial master prompt pre‑lock; after the autonomy lock the AI system autonomously performed literature review, study selection, preregistration, code generation, experiment execution, registered analysis, internal peer‑review checks, and manuscript drafting without manual human editing of technical content. Preregistration: CNSM2026‑AgenticNetOps‑GuardedGVR‑v1 (preregistration manifest sha256 recorded in the execution manifest). Orchestration and tools: hosted\_netops\_gvr\_v1 adapter, deterministic synthetic intent\ensuremath{\rightarrow}configuration task generator (fixed seeds), containerized emulator images (two topology classes), deterministic validator deterministic\_netops\_validator\_v5, and a preregistered generate--validate--repair policy allowing up to one repair. Execution summary (archived): planned\_episode\_count=400, completed\_episode\_count=374, failed\_episode\_count=26, model\_calls\_used=208, complete\_pairs=187. Analysis artifacts and recomputation: primary confirmatory analysis was produced by the registered executor 'paired\_binary\_analysis\_v1' and archived (analysis/paired\_contingency\_table.csv, analysis/condition\_summary.csv, analysis/missingness\_summary.csv, analysis/analysis\_log.jsonl). The preregistered sensitivity CI for the incomplete‑pair encoding was deterministically recomputed from archived contingency and missingness artifacts (analysis/paired\_contingency\_table.csv and analysis/missingness\_summary.csv) using bootstrap resamples=10000 and seed=1729; details of that recomputation are described in Methods/Results. Missing or unexecuted diagnostics: the registered analysis executor did not output the per‑episode paired latency diagnostics required to evaluate preregistered latency non‑inferiority H2; H2 remains unconfirmed. Post‑lock human editing: no manual human editing, rewriting, or selective improvement of technical content was performed after the autonomy lock. Retries and failures: handled per preregistration and logged in the execution manifest. Representative artifact pointers and their manifest entries (filenames and archived artifact references) are recorded in the archived execution manifest for deterministic verification and re‑execution of registered analysis contracts.

\begin{thebibliography}{99}
\bibitem{ref1} Muhammad Bilal, Jon Crowcroft, Ruizhi Wang, Xiaolong Xu, Schahram Dustdar, "Large Language Models for Agentic NetOps and AIOps: Architectures, Evaluation, and Safety", 2026
\bibitem{ref2} Andrew Brown, Muhammad Roman, Barry Devereux, "A Systematic Literature Review of Retrieval-Augmented Generation: Techniques, Metrics, and Challenges", 2025, 10.3390/bdcc9120320
\bibitem{ref3} PIerre Dantas, Lucas Cordeiro, Ehsan Nowroozi, Tihanyi Norbert, "Toward Safe LLM Agents: A Survey of Specification, Verification, and Enforcement", 2026, 10.48550/arxiv.2608.14590
\bibitem{ref4} Yue Zhang, Yafu Li, Leyang Cui, Cai Deng, Lemao Liu, Tingchen Fu, Xinting Huang, Enbo Zhao, Yanwen Zhang, Yulong Chen, Longyue Wang, Ahn Tuan Luu, Wei Bi, Freda Shi, Shuming Shi, "🧜Siren's Song in the AI Ocean: A Survey on Hallucination in Large Language Models", 2025, 10.1162/coli.a.16
\bibitem{ref5} Weikai Xu, Chengrui Huang, Shen Gao, Shuo Shang, "LLM-Based Agents for Tool Learning: A Survey", 2025, 10.1007/s41019-025-00296-9
\end{thebibliography}

\end{document}
